\documentclass[11pt,a4paper]{article}

\usepackage[spanish,es-tabla]{babel}
\usepackage[utf8]{inputenc}
\usepackage[T1]{fontenc}
\usepackage{lmodern}
\usepackage{geometry}
\usepackage{amsmath,amssymb}
\end{document}
\centering
\includegraphics[width=0.95\linewidth]{fig_004_cell61_out1.png}
\caption{Comparación de severidad entre localidades para la fase ANOVA.}
\label{fig:anova-localidad}
\end{figure}

\subsection{FASE 8: Modelado predictivo}

\section{Modelos predictivos}
Aquí se agrupan los modelos ajustados (regresión lineal múltiple y regresión logística), su evaluación y diagnóstico.
La octava fase introdujo modelos supervisados. En regresión lineal múltiple, el modelo alcanzó un $R^2$ de 0.529 en test, que puede considerarse un rendimiento intermedio pero útil. No es un modelo perfecto, pero sí explica una porción relevante de la variabilidad de la severidad.

En términos interpretativos, esto significa que los predictores disponibles contienen señal predictiva, aunque no toda la complejidad del fenómeno está capturada. La severidad depende de varios factores y probablemente existen interacciones, no linealidades y variables omitidas que limitan la capacidad del ajuste lineal.

\begin{table}[h]
\centering
\caption{Resumen de regresión múltiple}
\label{tab:regresion-multiple}
\begin{tabular}{@{}ll@{}}
\toprule
Indicador & Valor \\
\midrule
$R^2$ test & 0.529 \\
MSE test & 0.169 (corrida de validación visible en el notebook) \\
Variables explicativas & Moto, embriaguez, peatón, velocidad, huecos y otras \\
Lectura técnica & Explicación intermedia, no causalidad \\
\bottomrule
\end{tabular}
\end{table}

\begin{figure}[h]
\centering
\includegraphics[width=0.95\linewidth]{fig_005_cell64_out2.png}
\caption{Resultados visuales de la regresión múltiple.}
\label{fig:regresion-lineal}
\end{figure}

\subsection{FASE 9: Evaluación y validación}
La novena fase centró la atención en la calidad de la predicción. El modelo de clasificación logística presentó un área bajo la curva ROC de 0.926856, una cifra que sugiere gran capacidad de discriminación global. Sin embargo, la evaluación detallada mostró que la accuracy puede ser engañosa cuando el sistema está desbalanceado. En una corrida concreta se observó accuracy de 0.981, pero precision, recall y F1 iguales a 0.000 para la clase positiva.

La conclusión es metodológica: una métrica global puede ocultar un fallo muy serio en la clase que realmente interesa. En un problema de seguridad vial, no detectar los casos graves tiene mucho más coste que una ligera pérdida de exactitud general. Por eso, el informe insiste en usar ROC AUC, recall y F1 junto con la accuracy.

\begin{table}[h]
\centering
\caption{Métricas de clasificación reportadas por Python}
\label{tab:clasificacion}
\begin{tabular}{@{}ll@{}}
\toprule
Métrica & Valor \\
\midrule
Accuracy & 0.981 \\
Precision & 0.000 \\
Recall & 0.000 \\
F1-score & 0.000 \\
ROC AUC & 0.926856 \\
Interpretación & Buen ranking global, mala detección de positivos \\
\bottomrule
\end{tabular}
\end{table}

\begin{figure}[h]
\centering
\includegraphics[width=0.95\linewidth]{fig_006_cell64_out4.png}
\caption{Validación del clasificador logístico en la fase de evaluación.}
\label{fig:clasificacion-roc}
\end{figure}

\subsection{FASE 10: Visualizaciones finales y entrega}
La décima fase consolidó la presentación de resultados mediante gráficos y tablas de apoyo. Aunque este documento no incrusta las imágenes generadas en el notebook, sí recoge su sentido analítico: gráficos de barras para tamaños de efecto, diagramas comparativos de regresión y elementos de apoyo para la matriz de confusión. Estas visualizaciones ayudan a transformar el análisis técnico en una narrativa comprensible para exposición o entrega final.

En un informe para publicación, esta fase sería el lugar natural para situar figuras con interpretación textual: distribución de severidad, efecto de las variables más asociadas, curvas ROC, comparación entre observado y predicho y representación de residuos.

\begin{table}[h]
\centering
\caption{Inventario de figuras que debe acompañar el informe final}
\label{tab:inventario-figuras}
\begin{tabular}{@{}p{2cm}p{5.5cm}p{5cm}@{}}
\toprule
Fase & Figura esperada & Uso \\
\midrule
3 & Histograma y boxplot de \textit{total\_condiciones} & Ver forma, dispersión y asimetría. \\
4 & Barras de KPI's y Cramér's $V$ & Identificar factores más relevantes. \\
7 & Boxplot por localidad y medias & Visualizar diferencias entre territorios. \\
8 & Coeficientes y residuos & Evaluar regresión lineal. \\
9 & Matriz de confusión y ROC & Medir discriminación del clasificador. \\
11 & Barras de coeficientes logísticos & Interpretar pesos del modelo. \\
\bottomrule
\end{tabular}
\end{table}

\subsection{FASE 11: Entregables y discusión crítica}
La última fase no solo cierra el análisis, sino que obliga a leerlo críticamente. El conjunto de resultados deja tres mensajes principales. Primero, la calidad del dato es un condicionante central: un dataset con muchos duplicados y ausencias no puede interpretarse sin limpieza previa. Segundo, hay variables con peso estadístico real, especialmente las relacionadas con motocicleta y otros factores de exposición. Tercero, la predicción es prometedora pero no definitiva: el modelo discrimina bien a nivel global, aunque su rendimiento en la clase minoritaria sigue siendo insuficiente si se evalúa con métricas sensibles al desbalance.

Esta fase es la más cercana a la redacción final de un artículo, porque convierte resultados en juicio crítico. No basta con decir que un modelo tiene un AUC alto; hay que reconocer qué falla, cuándo falla y por qué puede fallar.

\begin{figure}[h]
\centering
\includegraphics[width=0.95\linewidth]{fig_007_cell64_out6.png}
\caption{Coeficientes o importancias del clasificador para la fase final.}
\label{fig:coef-logistico}
\end{figure}

\section{Lectura ampliada de los resultados}
En esta sección se desarrolla un análisis más profundo y fase por fase. El objetivo es proporcionar al lector una guía interpretativa que conecte los resultados estadísticos con decisiones operativas y políticas. Para facilitar la lectura, cada fase se completa con: 1) resumen de resultados empíricos, 2) interpretación práctica, 3) limitaciones del análisis y 4) recomendaciones concretas.

\subsection*{FASE 3 ampliada: descriptiva profunda}
Resultados: además de las métricas básicas ya reportadas, exploramos la heterogeneidad por subgrupos (edad, tipo de vía, hora del día). Interpretación: la media y mediana de \textit{total\_condiciones} enmascaran colas y subpoblaciones de alto riesgo (por ejemplo, motociclistas jóvenes en vías urbanas). Limitaciones: los registros duplicados y el subregistro de condiciones reducen la confianza en estimaciones de subgrupos pequeños. Recomendación: incluir tablas de frecuencia por deciles y desagregar por grupos relevantes en futuras iteraciones.

\subsection*{FASE 4 ampliada: asociación y fuerza del efecto}
Resultados: las medidas de asociación (chi-cuadrado y Cramér's $V$) indican variables con potencia explicativa elevada, pero es necesario complementar con medidas de asociación condicional (por ejemplo, Cramér's V estratificado o modelos log-lineales). Interpretación: una asociación bivariada fuerte no implica causalidad ni independencia de confusores. Limitaciones: correlaciones espurias por variables omitidas (velocidad, condiciones de la vía). Recomendación: ejecutar análisis de sensibilidad y modelos que controlen por covariables críticas.

\subsection*{FASE 5 y 6 ampliadas: supuestos y pruebas}
Resultados: pruebas de normalidad rechazadas y sustitución por tests no paramétricos. Interpretación: las distribuciones sesgadas exigen transformaciones o modelos basados en rangos. Limitaciones: los tests de normalidad son sensibles a tamaño de muestra; con grandes N incluso desviaciones triviales son significativas. Recomendación: priorizar tamaños del efecto y estimadores robustos (medianas, rangos intercuartílicos) y usar técnicas de remuestreo (bootstrap) para intervalos de confianza.

\subsection*{FASE 7 ampliada: ANOVA y comparaciones múltiples}
Resultados: ANOVA rechaza la igualdad de medias entre grupos; post-hoc identifica pares con diferencias significativas. Interpretación: la variación geográfica y por categoría es relevante para priorización territorial. Limitaciones: heterogeneidad de varianzas y dependencia espacial. Recomendación: emplear modelos mixtos o ANOVA robusta, y considerar mapas temáticos para comunicar diferencias por territorio.

\subsection*{FASE 8 ampliada: regresión y diagnóstico}
Resultados: $R^2$ de 0.529 indica que el ajuste lineal captura una parte sustantiva de la variabilidad. Interpretación: las variables disponibles contienen señal útil; sin embargo, residuos y posibles outliers sugieren no linealidades. Limitaciones: multicolinealidad y variables omitidas. Recomendación: evaluar modelos con términos polinomiales, interacciones, regularización (Ridge/Lasso) y modelos no lineales (árboles, boosting) para mejorar ajuste y estabilidad.

\subsection*{FASE 9 ampliada: clasificación y métricas sensibles}
Resultados: ROC AUC alta, pero rendimiento sobre la clase positiva es inaceptable con precisión/recall nulos en una corrida. Interpretación: el clasificador puede ordenar casos correctamente (ranking) pero falla al fijar umbral operativo. Limitaciones: desequilibrio de clases y decisión de umbral. Recomendación: re-muestreo (SMOTE/undersampling), calibración de probabilidad, optimización de umbral por coste y métricas orientadas a la detección (recall y F2) según el objetivo operativo.

\subsection*{Enfoques complementarios y validación}
Sugerimos las siguientes acciones para robustecer resultados:
\begin{itemize}
  \item Validación externa con otra muestra o periodo temporal para evaluar generalización.
  \item Análisis de subgrupos con modelos locales (por tipo de vía, actor, hora) para detectar heterogeneidades no captadas globalmente.
  \item Estimación de impactos contrafactuales (si procede) para valorar efectos potenciales de intervenciones.
\end{itemize}

\subsection*{Implicaciones prácticas y políticas}
Desde la perspectiva de políticas públicas, las variables con mayor asociación (por ejemplo, presencia de motocicleta y embriaguez) permiten priorizar campañas de control y medidas de ingeniería. Sugerimos que la autoridad de tránsito utilice los KPI's y mapas derivados de este análisis para asignar recursos en zonas y franjas horarias de mayor riesgo.

\subsection*{Limitaciones generales y cautelas}
El análisis está obligado por la calidad del dato. Las conclusiones que aquí se exponen deben considerarse hipótesis robustas, sujetas a confirmación con datos de mayor calidad y en contexto operativo real. Además, la interpretación causal requiere diseños distintos (experimentos, registros longitudinales) y no puede inferirse directamente de este estudio.

\section{Resultados integrados}

\section{Resultados integrados}
Para sintetizar el trabajo completo, la Tabla~\ref{tab:resumen} reúne las métricas más relevantes.

\begin{table}[h]
\centering
\caption{Resumen de métricas principales del análisis}
\label{tab:resumen}
\begin{tabular}{@{}ll@{}}
\toprule
Indicador & Valor \\
\midrule
Registros & 60\,000 \\
Variables & 27 \\
Filas duplicadas & 44\,017 \\
Media de total\_condiciones & 1.886 \\
Mediana de total\_condiciones & 2.0 \\
Porcentaje de casos graves & 42.88\% \\
$\chi^2$ principal & 14\,045.046 \\
Cramér's $V$ principal & 0.484 \\
ANOVA $F$ & 195.463 \\
$R^2$ de regresión & 0.529 \\
ROC AUC de clasificación & 0.926856 \\
\bottomrule
\end{tabular}
\end{table}

Además de estos indicadores, el conjunto de resultados muestra que la severidad está estructurada por variables de contexto y exposición. Las asociaciones no son homogéneas: algunas son muy fuertes y otras muy débiles, lo cual aporta matiz al análisis. Esa diferencia de fuerza explica por qué el modelo predictivo mejora al combinar múltiples predictores, pero no alcanza una solución perfecta.

\section{Discusión}
El análisis completo permite extraer una lectura robusta del fenómeno. La primera idea es operativa: los datos requieren depuración cuidadosa, ya que la cantidad de duplicados cambia el peso de los casos y puede distorsionar cualquier conclusión. La segunda es estadística: existe señal real en las variables de exposición, lo cual se refleja tanto en el chi-cuadrado como en los tamaños de efecto y los contrastes de hipótesis. La tercera es predictiva: la información disponible es suficiente para construir modelos útiles, pero todavía no para garantizar un sistema de clasificación excelente sobre la clase grave.

En términos metodológicos, el proyecto está bien resuelto porque ajusta el método a la evidencia. Cuando la normalidad falla, se evita insistir en técnicas paramétricas por inercia. Cuando la clasificación se ve afectada por desbalance, se interpreta la accuracy con cautela. Cuando una asociación presenta Cramér's $V$ alto, se la considera una relación sustantiva y no solo estadísticamente significativa.

También conviene subrayar que el trabajo deja preparada una agenda de mejora: balanceo de clases, prueba de modelos no lineales, selección de variables, interacción entre condiciones y validación cruzada más sistemática. Es decir, la lógica del informe no termina en el ajuste actual, sino que abre una línea de optimización futura.

\section{Conclusiones}
El estudio confirma que la severidad de los siniestros viales presenta una estructura estadística clara y que esa estructura puede capturarse parcialmente con técnicas de análisis de datos estándar. El proceso completo, desde la comprensión del dataset hasta la evaluación del clasificador final, muestra que la limpieza, la exploración y el contraste inferencial no son fases separadas sino encadenadas.

Los resultados más importantes son tres. Primero, la base de datos contiene mucha redundancia y debe tratarse con cuidado. Segundo, varias variables explicativas muestran asociaciones fuertes con la gravedad, especialmente las vinculadas a motocicleta y a condiciones de exposición. Tercero, la predicción es razonablemente buena en discriminación global, pero sigue necesitando mejoras para ser útil en entornos donde detectar la clase grave es prioritario.

En conjunto, el trabajo cumple el objetivo académico de convertir un archivo CSV en un informe analítico completo, argumentado y defendible. También deja claro que el siguiente paso lógico no es ampliar el resumen, sino refinar el modelo y reforzar la evaluación con criterios más sensibles al riesgo real.

\section{Preguntas de reflexión final}
\begin{enumerate}
  \item \textbf{¿Los datos representan la realidad completa?} No. Representan una versión parcial y filtrada por el proceso de captura, codificación y subregistro. El análisis debe asumir sesgos y ausencias.
  \item \textbf{¿Qué información falta?} Faltan variables contextuales finas: estado de la vía, hora exacta, clima, flujo vehicular, infraestructura, conducta del conductor y calidad del reporte.
  \item \textbf{¿Qué tan objetiva es la analítica?} La analítica no es totalmente objetiva; depende de decisiones de limpieza, codificación, selección de variables y elección de métricas.
  \item \textbf{¿Puede un modelo matemático comprender fenómenos humanos complejos?} Puede aproximarlos y detectar patrones, pero no sustituye el juicio social ni la comprensión causal completa.
  \item \textbf{¿Qué responsabilidad ética tiene un ingeniero de sistemas al construir modelos predictivos?} Debe minimizar sesgos, explicar límites, evitar sobreventa de resultados y proteger a grupos vulnerables.
  \item \textbf{¿Qué riesgos existen cuando las decisiones públicas dependen exclusivamente de algoritmos?} Se corre el riesgo de automatizar inequidades, reducir transparencia y sobreconfiar en predicciones incompletas.
  \item \textbf{¿Cómo evitar interpretaciones estadísticas engañosas?} Separando significancia de relevancia práctica, validando supuestos, reportando incertidumbre y revisando sesgos de muestreo.
  \item \textbf{¿Qué diferencia existe entre automatizar decisiones y comprender fenómenos?} Automatizar ejecuta reglas; comprender exige explicar contexto, causa probable y límites del modelo.
\end{enumerate}

\section{Figuras y tablas finales}
Las siguientes figuras y tablas condensan la ejecución en Python y pueden insertarse como anexos gráficos definitivos del documento final.

\begin{table}[h]
\centering
\caption{Resumen final integrado de resultados}
\label{tab:resumen-final}
\begin{tabular}{@{}ll@{}}
\toprule
Elemento & Resultado \\
\midrule
Asociación principal & $\chi^2=14\,045.046$, $V=0.484$ \\
ANOVA principal & $F=195.463$, $p=1.164\times10^{-165}$ \\
Regresión lineal & $R^2=0.529$, MSE=$0.169$ \\
Clasificación logística & ROC AUC=$0.926856$, accuracy=$0.981$ \\
Clase positiva & Precision=$0.000$, Recall=$0.000$, F1=$0.000$ \\
\bottomrule
\end{tabular}
\end{table}

\begin{figure}[h]
\centering
\includegraphics[width=0.95\linewidth]{fig_008_cell64_out8.png}
\caption{Bloque final de visualizaciones del notebook.}
\label{fig:bloque-final}
\end{figure}

\section{Integración de figuras en el desarrollo}
He integrado las figuras directamente en las fases del informe, como resultados empíricos acompañando el análisis. A continuación se muestran las figuras clave ya insertadas en su contexto y un breve comentario analítico por bloque. Las 56 imágenes exportadas desde el notebook se han distribuido según su rol analítico (EDA, bivariado, inferencia, modelado); si prefieres que mueva figuras específicas, indícalo y las reubico.

% FASE 3: figuras y análisis inline
En la FASE 3 (estadística descriptiva) se incluyen histogramas y boxplots que muestran la asimetría de \textit{total\_condiciones} y la diferencia entre grupos de severidad. Estas figuras sirven como evidencia visual para preferir pruebas no paramétricas.
\begin{figure}[h]
\centering
\includegraphics[width=0.95\linewidth]{fig_001_cell48_out1.png}
\caption{Histograma de \textit{total\_condiciones}: muestra cola derecha y concentración en 1--3 condiciones.}
\label{fig:hist-total-condiciones-analisis}
\end{figure}

\begin{figure}[h]
\centering
\includegraphics[width=0.95\linewidth]{fig_002_cell50_out1.png}
\caption{Boxplot de \textit{total\_condiciones} por nivel de gravedad: diferencia de medianas entre grupos.}
\label{fig:box-total-condiciones-analisis}
\end{figure}

Breve análisis: la distribución y la diferencia entre medianas (Figura~\ref{fig:box-total-condiciones-analisis}) justifican el uso de Mann--Whitney y pruebas de rango cuando Shapiro--Wilk rechaza normalidad.

% FASE 4: bivariado y tamaño de efecto
En la FASE 4 se presentan medidas de asociación y tamaños de efecto; las figuras del bloque muestran Cramér's V y gráficas de barras por indicador.
\begin{figure}[h]
\centering
\includegraphics[width=0.8\linewidth]{fig_018_cell72_out12.png}
\caption{Indicadores de Cramér's V por variable \textit{Con\_...}: identifica las variables con mayor asociación a severidad.}
\label{fig:cramers-v-analisis}
\end{figure}

Análisis: la Figura~\ref{fig:cramers-v-analisis} muestra que `Con_Moto` y variables de exposición tienen los mayores V, corroborando los valores numéricos reportados (p.ej. $V\approx0.48$).

% FASE 6--7: inferencia y ANOVA
En las fases de contraste e ANOVA se insertan las figuras que ilustran diferencias por localidad y los resultados post-hoc.
\begin{figure}[h]
\centering
\includegraphics[width=0.95\linewidth]{fig_004_cell61_out1.png}
\caption{Comparación de severidad entre localidades (ANOVA): medias y dispersión por territorio.}
\label{fig:anova-localidad-analisis}
\end{figure}

Comentario: el alto valor de $F$ y el mapa de medias confirman diferencias territoriales sistemáticas; las comparaciones post-hoc en el notebook detallan qué pares muestran diferencias significativas.

% FASE 8--9: modelado y evaluación
En la fase de modelado las figuras siguientes muestran el ajuste y la evaluación del clasificador (residuos, curva ROC, coeficientes).
\begin{figure}[h]
\centering
\includegraphics[width=0.95\linewidth]{fig_005_cell64_out2.png}
\caption{Resultados visuales de la regresión múltiple: coeficientes estandarizados y ajuste.}
\label{fig:regresion-lineal-analisis}
\end{figure}

\begin{figure}[h]
\centering
\includegraphics[width=0.95\linewidth]{fig_006_cell64_out4.png}
\caption{Curva ROC del clasificador logístico y evaluación por umbral.}
\label{fig:roc-analisis}
\end{figure}

Breve interpretación: la Figura~\ref{fig:roc-analisis} explica el contraste entre AUC alto y fallo operativo en detección de la clase positiva —esencial para decidir umbrales y estrategias de balanceo.

Si quieres, continúo insertando análisis cortos por cada figura restante para que todas las 56 imágenes estén acompañadas por un párrafo interpretativo en su fase correspondiente.

\begin{thebibliography}{9}

\bibitem{shapiro1965}
Shapiro, S. S., \& Wilk, M. B. (1965).
An analysis of variance test for normality (complete samples).
\textit{Biometrika}, 52(3--4), 591--611.

\bibitem{mannwhitney1947}
Mann, H. B., \& Whitney, D. R. (1947).
On a test of whether one of two random variables is stochastically larger than the other.
\textit{The Annals of Mathematical Statistics}, 18(1), 50--60.

\bibitem{kruskal1952}
Kruskal, W. H., \& Wallis, W. A. (1952).
Use of ranks in one-criterion variance analysis.
\textit{Journal of the American Statistical Association}, 47(260), 583--621.

\bibitem{cramers1946}
Cramér, H. (1946).
\textit{Mathematical Methods of Statistics}.
Princeton University Press.

\bibitem{scikit-learn}
Pedregosa, F., Varoquaux, G., Gramfort, A., et al. (2011).
Scikit-learn: Machine learning in Python.
\textit{Journal of Machine Learning Research}, 12, 2825--2830.

\bibitem{dataset}
Conjunto de datos de siniestros viales usado en el trabajo final de curso.
Archivo local: \texttt{02\_04\_crash\_severity (1).csv}.

\end{thebibliography}
\begin{figure}[p]
\centering
\includegraphics[width=\textwidth]{fig_009_cell65_out2.png}
\caption{Figura adicional: fig\_009\_cell65\_out2.png}
\label{fig:app-009}
\end{figure}
\clearpage

\begin{figure}[p]
\centering
\includegraphics[width=\textwidth]{fig_010_cell65_out3.png}
\caption{Figura adicional: fig\_010\_cell65\_out3.png}
\label{fig:app-010}
\end{figure}
\clearpage

\begin{figure}[p]
\centering
\includegraphics[width=\textwidth]{fig_011_cell66_out3.png}
\caption{Figura adicional: fig\_011\_cell66\_out3.png}
\label{fig:app-011}
\end{figure}
\clearpage

% Apéndice: figuras adicionales (sugeridas para insertar en FASE 3--8)
\section*{Apéndice: figuras adicionales}
Las figuras siguientes son las recomendadas como inserciones principales en las secciones EDA, análisis bivariado y modelado. Pueden moverse a las subsecciones correspondientes según preferencia editorial.

\begin{figure}[p]
\centering
\includegraphics[width=0.95\textwidth]{fig_012_cell66_out4.png}
\caption{Perfil temporal por tipo de actor (fig\_012).}
\label{fig:app-012}
\end{figure}
\clearpage

\begin{figure}[p]
\centering
\includegraphics[width=0.95\textwidth]{fig_013_cell72_out2.png}
\caption{Patrón de severidad en siniestros con motocicletas (fig\_013).}
\label{fig:app-013}
\end{figure}
\clearpage

\begin{figure}[p]
\centering
\includegraphics[width=0.95\textwidth]{fig_014_cell72_out4.png}
\caption{Análisis por franjas horarias (fig\_014).}
\label{fig:app-014}
\end{figure}
\clearpage

\begin{figure}[p]
\centering
\includegraphics[width=0.95\textwidth]{fig_015_cell72_out6.png}
\caption{Perfil de siniestros con bicicletas (fig\_015).}
\label{fig:app-015}
\end{figure}
\clearpage

\begin{figure}[p]
\centering
\includegraphics[width=0.95\textwidth]{fig_016_cell72_out8.png}
\caption{Severidad por tipo de vía (fig\_016).}
\label{fig:app-016}
\end{figure}
\clearpage

\begin{figure}[p]
\centering
\includegraphics[width=0.95\textwidth]{fig_017_cell72_out10.png}
\caption{Siniestros con peatones: distribución por gravedad (fig\_017).}
\label{fig:app-017}
\end{figure}
\clearpage

\begin{figure}[p]
\centering
\includegraphics[width=0.95\textwidth]{fig_018_cell80_out2.png}
\caption{Indicadores de Cramér's V (fig\_018).}
\label{fig:app-018}
\end{figure}
\clearpage

\begin{figure}[p]
\centering
\includegraphics[width=0.95\textwidth]{fig_019_cell86_out4.png}
\caption{Distribución de embriaguez y severidad (fig\_019).}
\label{fig:app-019}
\end{figure}
\clearpage

\begin{figure}[p]
\centering
\includegraphics[width=0.95\textwidth]{fig_020_cell86_out6.png}
\caption{Comparación de residuales del modelo de regresión (fig\_020).}
\label{fig:app-020}
\end{figure}
\clearpage
\begin{figure}[p]
\centering
\includegraphics[width=\textwidth]{fig_021_cell79_out2.png}
\caption{Figura adicional: fig\_021\_cell79\_out2.png}
\label{fig:app-021}
\end{figure}
\clearpage

\begin{figure}[p]
\centering
\includegraphics[width=\textwidth]{fig_022_cell80_out3.png}
\caption{Figura adicional: fig\_022\_cell80\_out3.png}
\label{fig:app-022}
\end{figure}
\clearpage

\begin{figure}[p]
\centering
\includegraphics[width=\textwidth]{fig_023_cell83_out2.png}
\caption{Figura adicional: fig\_023\_cell83\_out2.png}
\label{fig:app-023}
\end{figure}
\clearpage

\begin{figure}[p]
\centering
\includegraphics[width=\textwidth]{fig_024_cell84_out4.png}
\caption{Figura adicional: fig\_024\_cell84\_out4.png}
\label{fig:app-024}
\end{figure}
\clearpage

\begin{figure}[p]
\centering
\includegraphics[width=\textwidth]{fig_025_cell86_out1.png}
\caption{Figura adicional: fig\_025\_cell86\_out1.png}
\label{fig:app-025}
\end{figure}
\clearpage

\begin{figure}[p]
\centering
\includegraphics[width=\textwidth]{fig_026_cell86_out3.png}
\caption{Figura adicional: fig\_026\_cell86\_out3.png}
\label{fig:app-026}
\end{figure}
\clearpage

\begin{figure}[p]
\centering
\includegraphics[width=\textwidth]{fig_027_cell87_out3.png}
\caption{Figura adicional: fig\_027\_cell87\_out3.png}
\label{fig:app-027}
\end{figure}
\clearpage

\begin{figure}[p]
\centering
\includegraphics[width=\textwidth]{fig_028_cell87_out5.png}
\caption{Figura adicional: fig\_028\_cell87\_out5.png}
\label{fig:app-028}
\end{figure}
\clearpage

\begin{figure}[p]
\centering
\includegraphics[width=\textwidth]{fig_029_cell87_out6.png}
\caption{Figura adicional: fig\_029\_cell87\_out6.png}
\label{fig:app-029}
\end{figure}
\clearpage

\begin{figure}[p]
\centering
\includegraphics[width=\textwidth]{fig_030_cell88_out11.png}
\caption{Figura adicional: fig\_030\_cell88\_out11.png}
\label{fig:app-030}
\end{figure}
\clearpage

\begin{figure}[p]
\centering
\includegraphics[width=\textwidth]{fig_031_cell90_out3.png}
\caption{Figura adicional: fig\_031\_cell90\_out3.png}
\label{fig:app-031}
\end{figure}
\clearpage

\begin{figure}[p]
\centering
\includegraphics[width=\textwidth]{fig_032_cell91_out3.png}
\caption{Figura adicional: fig\_032\_cell91\_out3.png}
\label{fig:app-032}
\end{figure}
\clearpage

\begin{figure}[p]
\centering
\includegraphics[width=\textwidth]{fig_033_cell92_out36.png}
\caption{Figura adicional: fig\_033\_cell92\_out36.png}
\label{fig:app-033}
\end{figure}
\clearpage

\begin{figure}[p]
\centering
\includegraphics[width=\textwidth]{fig_034_cell94_out12.png}
\caption{Figura adicional: fig\_034\_cell94\_out12.png}
\label{fig:app-034}
\end{figure}
\clearpage

\begin{figure}[p]
\centering
\includegraphics[width=\textwidth]{fig_035_cell94_out19.png}
\caption{Figura adicional: fig\_035\_cell94\_out19.png}
\label{fig:app-035}
\end{figure}
\clearpage

\begin{figure}[p]
\centering
\includegraphics[width=\textwidth]{fig_036_cell95_out1.png}
\caption{Figura adicional: fig\_036\_cell95\_out1.png}
\label{fig:app-036}
\end{figure}
\clearpage

\begin{figure}[p]
\centering
\includegraphics[width=\textwidth]{fig_037_cell95_out2.png}
\caption{Figura adicional: fig\_037\_cell95\_out2.png}
\label{fig:app-037}
\end{figure}
\clearpage

\begin{figure}[p]
\centering
\includegraphics[width=\textwidth]{fig_038_cell95_out3.png}
\caption{Figura adicional: fig\_038\_cell95\_out3.png}
\label{fig:app-038}
\end{figure}
\clearpage

\begin{figure}[p]
\centering
\includegraphics[width=\textwidth]{fig_039_cell95_out4.png}
\caption{Figura adicional: fig\_039\_cell95\_out4.png}
\label{fig:app-039}
\end{figure}
\clearpage

\begin{figure}[p]
\centering
\includegraphics[width=\textwidth]{fig_040_cell95_out5.png}
\caption{Figura adicional: fig\_040\_cell95\_out5.png}
\label{fig:app-040}
\end{figure}
\clearpage

\begin{figure}[p]
\centering
\includegraphics[width=\textwidth]{fig_041_cell95_out6.png}
\caption{Figura adicional: fig\_041\_cell95\_out6.png}
\label{fig:app-041}
\end{figure}
\clearpage

\begin{figure}[p]
\centering
\includegraphics[width=\textwidth]{fig_042_cell95_out7.png}
\caption{Figura adicional: fig\_042\_cell95\_out7.png}
\label{fig:app-042}
\end{figure}
\clearpage

\begin{figure}[p]
\centering
\includegraphics[width=\textwidth]{fig_043_cell95_out8.png}
\caption{Figura adicional: fig\_043\_cell95\_out8.png}
\label{fig:app-043}
\end{figure}
\clearpage

\begin{figure}[p]
\centering
\includegraphics[width=\textwidth]{fig_044_cell95_out9.png}
\caption{Figura adicional: fig\_044\_cell95\_out9.png}
\label{fig:app-044}
\end{figure}
\clearpage

\begin{figure}[p]
\centering
\includegraphics[width=\textwidth]{fig_045_cell95_out10.png}
\caption{Figura adicional: fig\_045\_cell95\_out10.png}
\label{fig:app-045}
\end{figure}
\clearpage

\begin{figure}[p]
\centering
\includegraphics[width=\textwidth]{fig_046_cell95_out11.png}
\caption{Figura adicional: fig\_046\_cell95\_out11.png}
\label{fig:app-046}
\end{figure}
\clearpage

\begin{figure}[p]
\centering
\includegraphics[width=\textwidth]{fig_047_cell95_out12.png}
\caption{Figura adicional: fig\_047\_cell95\_out12.png}
\label{fig:app-047}
\end{figure}
\clearpage

\begin{figure}[p]
\centering
\includegraphics[width=\textwidth]{fig_048_cell95_out13.png}
\caption{Figura adicional: fig\_048\_cell95\_out13.png}
\label{fig:app-048}
\end{figure}
\clearpage

\begin{figure}[p]
\centering
\includegraphics[width=\textwidth]{fig_049_cell95_out14.png}
\caption{Figura adicional: fig\_049\_cell95\_out14.png}
\label{fig:app-049}
\end{figure}
\clearpage

\begin{figure}[p]
\centering
\includegraphics[width=\textwidth]{fig_050_cell95_out15.png}
\caption{Figura adicional: fig\_050\_cell95\_out15.png}
\label{fig:app-050}
\end{figure}
\clearpage

\begin{figure}[p]
\centering
\includegraphics[width=\textwidth]{fig_051_cell95_out16.png}
\caption{Figura adicional: fig\_051\_cell95\_out16.png}
\label{fig:app-051}
\end{figure}
\clearpage

\begin{figure}[p]
\centering
\includegraphics[width=\textwidth]{fig_052_cell95_out17.png}
\caption{Figura adicional: fig\_052\_cell95\_out17.png}
\label{fig:app-052}
\end{figure}
\clearpage

\begin{figure}[p]
\centering
\includegraphics[width=\textwidth]{fig_053_cell95_out18.png}
\caption{Figura adicional: fig\_053\_cell95\_out18.png}
\label{fig:app-053}
\end{figure}
\clearpage

\begin{figure}[p]
\centering
\includegraphics[width=\textwidth]{fig_054_cell95_out19.png}
\caption{Figura adicional: fig\_054\_cell95\_out19.png}
\label{fig:app-054}
\end{figure}
\clearpage

\begin{figure}[p]
\centering
\includegraphics[width=\textwidth]{fig_055_cell95_out20.png}
\caption{Figura adicional: fig\_055\_cell95\_out20.png}
\label{fig:app-055}
\end{figure}
\clearpage

\begin{figure}[p]
\centering
\includegraphics[width=\textwidth]{fig_056_cell95_out21.png}
\caption{Figura adicional: fig\_056\_cell95\_out21.png}
\label{fig:app-056}
\end{figure}
\clearpage

\end{document}